\cvsection{Expériences professionnelles}

\begin{cventries}
  \cventry
    {Stagiaire ingénieur de recherche en IA}
    {INRIA}
    {Bordeaux, France}
    {Mars 2026 -- Septembre 2026}
    {
      \begin{cvitems}
        \item {Conception et développement en autonomie d’une application et d’un agent IA multi-étapes pour aider les chercheurs à trouver des articles scientifiques.}
        \item {Construction de l’agent en Python sans framework d’orchestration, avec API OpenAI, prompts, appels d’outils et gestion de la mémoire et du contexte.}
        \item {Développement d’un pipeline RAG avec chunking, embeddings et recherche sémantique dans ChromaDB, associé à une persistance MongoDB.}
        \item {Développement d’API REST Python/FastAPI documentées avec OpenAPI et d’une interface React/Next.js pour l’utilisation de l’agent.}
        \item {Automatisation de la collecte de métadonnées d’articles scientifiques en ligne et transmission des résultats de tâches cron par webhooks.}
        \item {Conteneurisation des composants avec Docker, versionnement sous Git et mise en place d’une CI/CD sur GitHub.}
        \item {Déploiement et serving de LLM avec vLLM sur une infrastructure HPC équipée de GPU NVIDIA et utilisation de Slurm pour les expérimentations.}
        \item {Évaluation de Qwen, Gemma, GLM et de plusieurs quantifications selon la qualité, la latence et le contexte, avec rédaction en cours d’un article scientifique.}
        \item {Gestion d’une application utilisant des modèles locaux avec LM Studio sur Mac Studio, avec installation et configuration des modèles.}
        \item {Évaluation d’OpenClaw dans un environnement Docker isolé et analyse des risques liés à l’accès autonome aux outils et aux fichiers.}
      \end{cvitems}
    }

  \cventry
    {Stagiaire de recherche en apprentissage par renforcement}
    {Université Karunya}
    {Coimbatore, Inde}
    {Juin 2025 -- Sept. 2025}
    {
      \begin{cvitems}
        \item {Développement d’un modèle d’apprentissage par renforcement et de son environnement d’entraînement pour la navigation autonome d’un robot hospitalier.}
        \item {Déploiement et optimisation du logiciel et du modèle sur Raspberry Pi pour des tests en conditions réelles.}
      \end{cvitems}
    }
\end{cventries}
